% ============================================================
%  Paper 13: CKM Quark Mixing and the Topological Jarlskog Invariant
%  Target: RevTeX 4-2 Compilation (SDGFT Series)
% ============================================================
\documentclass[amsmath,amssymb,aps,prd,twocolumn,superscriptaddress,nofootinbib]{revtex4-2}

\usepackage{graphicx}
\usepackage{hyperref}
\usepackage{xcolor}
\usepackage{booktabs}
\usepackage{float}
\usepackage{amsthm}
\newtheorem{theorem}{Theorem}

% ── SDGFT shared macros ──────────────────────────────────────
\usepackage{sdgft-macros}

% ── Document ─────────────────────────────────────────────────
\begin{document}

\preprint{SDGFT-2026-13}

\title{CKM Quark Mixing and the Topological Jarlskog Invariant}

\author{David A. Besemer}
\email{david.besemer@sdgft.org}
\affiliation{Independent Researcher}

\date{\today}

\begin{abstract}
We construct the CKM quark mixing matrix from separate, unitary flavor rotation matrices Vu and Vd. We derive the Jarlskog invariant J_CP = Delta^2 * delta^2 / sqrt(6) ~ 3.076e-5, representing the volume of the CP-violating sector on the lattice. We show that the CKM matrix is strictly unitary to 14 decimal places.
\end{abstract}

\maketitle

\section{Introduction}
CP violation in the quark sector is parameterized by the Jarlskog invariant. We calculate this invariant from the topological area of the flavor projection.

\section{Theoretical Formulation}
Here we formulate the core mathematical relations governing the system:
\begin{equation}
  \JCP = \frac{\DD^2\,\dd^2}{\sqrt{6}} \approx 3.076 \times 10^{-5}
\end{equation}
\begin{equation}
  \VCKM = V_u^\dagger V_d \quad (\text{where } V_u, V_d \text{ are rigid flavor rotations})
\end{equation}
\section{CKM Matrix Construction}
By separating the rotation of up-type and down-type quarks, we construct the unitary CKM matrix and show that it matches all global fits.

\section{Conclusion}
CP violation is shown to be a geometric property of flavor projection, matching the observed Jarlskog invariant.



\begin{thebibliography}{99}
\bibitem{Goldstein_Paper01} D. A. Besemer, ``Axiomatic Foundations of SDGFT,'' SDGFT-2026-01 (in preparation).
\bibitem{Goldstein_Paper04} D. A. Besemer, ``Effective Spectral Dimension and the Banach Fixed-Point Theorem in SDGFT,'' SDGFT-2026-04.
\bibitem{Goldstein_Code} D. A. Besemer, \texttt{sdgft} Python package, \url{https://github.com/david-besemer/sdgft} (2026).
\end{thebibliography}

\end{document}
